% Resumen breve para la contraportada (revisión C4, 23/11/2026). El texto completo de
% resumentexto.tex no cabe en el bloque de color con la letra \large; se acorta el texto en
% lugar de reducir la letra. Si cambia un resultado en resumentexto.tex, revisar también este.
\vspace{0.7cm}

Este trabajo presenta el diseño e implementación de VeriVoto, un prototipo SaaS
multi-tenant de votación digital para organizaciones: gestión de procesos, acceso
del votante, emisión del voto, recibos, auditoría y escrutinio. La arquitectura
integra cinco servicios propios, Keycloak, PostgreSQL y comunicación REST con
patrón \emph{outbox}. La evaluación, en local y con datos sintéticos, aceptó la
accesibilidad del flujo del votante (H1a), con auditoría automática y revisión
manual con lector de pantalla, y la usabilidad con cinco votantes (H1b) y con
cinco administradores sin formación previa (RNF-4). La prueba de carga rechazó
H1c: con 500 usuarios simultáneos no se alcanzaron los dos segundos. El trabajo
describe el prototipo y sus límites, sin atribuirle resistencia a la coacción ni
verificabilidad extremo a extremo.

\vspace{0.7cm}

\section*{Abstract}
This work presents the design and implementation of VeriVoto, a multi-tenant
voting Software as a Service prototype for organizations: election management,
voter access, ballot submission, receipts, audit records and tallying. The
architecture integrates five application services, Keycloak, PostgreSQL and REST
communication with an outbox pattern. The evaluation, run locally with synthetic
data, accepted voter-flow accessibility (H1a), with an automated audit and a
manual screen reader review, and usability with five voters (H1b) and with five
administrators without previous training (RNF-4). The load test rejected H1c: with
500 simultaneous users, the two-second threshold was not met. This report
describes the prototype and its limitations and does not claim coercion
resistance or end-to-end verifiability.
